\documentclass[../main.tex]{subfiles}

\begin{document}

\section{Introduction to Large Language Models}

You probably have heard of Large Language Models (LLMs)
and even have your favorite ones. You may have also noticed
that if you and your friend ask a same question to a LLM,
you get different response, or at least outputs that are
not identical from word to word. That is because LLMs
are generative models. As always, let's get started by
getting used to few vocabs.

\begin{definition}{}{}
\textbf{Generative models} are models that predicts and
estimates the values on a probability distribution.
\end{definition}

A generative model is also capable of taking in sequential
data for a more context-rich prediction. An example of
generative models would be an autoregressive model.

\begin{definition}{}{}
An \textbf{autoregressive model} is a generative model
that predicts the future values in a sequence given the
past input sequence.
\end{definition}

For instance, when given the sequence of integers
$0, 1, 1, 2, 3, 5, 8$, an autoregressive model is likely
to generate $13$ by using previously learned pattern and
the input. In other words, it is simply a model that
predicts the term $P(x_n \mid x_1, \ldots x_{n-1})$.

\begin{definition}{}{}
A \textbf{language model} is a model that utilize human
languages and predicts the next tokens given the previous
inputs.
\end{definition}

An interesting fact about such models is that they sometimes
exhibit \textbf{emergent behavior} where they actually
``learn'' and exhibit capabilities that were not explicitly
present in the training data.

\subsection{Tokenization}

We know that at the end, language models are huge machines
that predict the next numerical values. Then how do such
models generate outputs in human language? This is where
\textbf{tokenization} kicks in.

\begin{definition}{}{}
\textbf{Tokenization} is a process of breaking down inputs
into smaller units called \textbf{tokens} which are often
represented as numeric values.
\end{definition}

To tokenize inputs, a tool called \textbf{tokenizer} is
used with \textbf{Byte Pair Encoding (BPE) algorithm}.
Intuitively, what we could do is assign numbers to each
alphabets like $1$ for $a$, $2$ for $b$, and so on. However,
this would be very inefficient considering the frequency of
each letter and scalability. Instead, what we can do with
the BPE algorithm is create a byte-efficient list of
merged chunks (that are not necessarily a single letter)
based on frequency. Here is an example from the
\href{https://platform.openai.com/tokenizer}{OpenAI Tokenizer
for GPT-5.x \& O1/3}.

Consider the following sentence:
\[
\text{Hippopotomonstrosesquippedaliophobia means fear of long words.}
\]
When tokenized, this $62$ characters sentence turns into $17$
tokens as following.
\begin{align*}
    &\text{[`H', `ipp', `opot', `omon', `st', `ros', `es', `qu',
            `ipped', `ali', `ophobia',} \\
    &\qquad\qquad\qquad\qquad\qquad
    \text{` means', ` fear', ` of', ` long',
          ` words', `.']}
\end{align*}
The token IDs of the tokens above are the following numbers.
\begin{align*}
    \text{[39, 3012, 101543, 43960, 302, 2199, 268, }
    &\text{351, 8193, 3010, 141868,} \\
    &\text{4748, 11747, 328, 1701, 6391, 13]}
\end{align*}
Now, to help the models ``understand'' the tokens and take
the IDs as input, we generally turn them into vectors through
the process of \textbf{One-Hot Encoding (OHE)}, though the process
may differ in modern Deep Learning for efficiency.

With OHE, the token ID $i$ can be translated to a vector with
$1$ in index $i$ and $0$ in other slots. However, notice that
OHE itself is not highly effective. To make the tokenizer
more ``meaningful'' and ``effective'', we could use the process
of embeddings. \textbf{Learned embedding vectors} uses somewhat
low dimensional vectors (usually a few thousand dimensions) to
make the vectors ``meaningful''. For example, in learned
embedding vector, the difference between the vectors that
represent king and queen will be similar to the difference
between the vectors that represent man and woman. Modern models
use \textbf{relative positional embeddings} to effectively
distinguish the position of the words. (Let's save this topic
for latter notes on transformers.)

\subsection{Language Model Output}

As you probably have guessed, language models do not simply
output the texts that we see. An autoregressive transformer
language model will output a sequence of multi-dimensional
vectors with probability distribution for the next token based
on the sequence of multi-dimension input vectors.

\begin{definition}{}{}
\textbf{Logits} are the \textbf{scores} of the next possible
token represented by individual vectors in the output sequence.
\end{definition}

Note that logits are \textit{raw} scores. To convert them to
probability distribution, we use something called
\textbf{softmax}, in which we will discuss further in notes
about neural networks.

\begin{definition}{}{}
\textbf{Decoding} is the process selecting the next token to
convert raw and abstract model output into a human language
with \textbf{detokenization}.
\end{definition}

There are different strategies in decoding, but the two main
methods are \textbf{greedy decoding} and \textbf{sampling}.
In greedy decoding, you select the token with the highest
probability. Meaning, the temperature (we will cover this
is neural nets notes) does not influence the selection.
Unlike greedy decoding, temperature does influence the result
from sampling. In sampling, you randomly select the next token
with accounting the individual probability. To put it simply,
imagine you have a biased die with thousands of sides. You
roll the die with the bias in each side, and the next token
is selected by the result.

Model outputs are random and highly organized. Take a look
at the following example. In the first case, say the model
input is ``I love'' and the expected model output is ``I
love potatoes''. In the second case, the model input is
``Veni, vidi,'' and the expected output is ``Veni, vidi,
vici''. When the model returns an output, which case do you
think the model has the higher probability of returning the
expected output? To answer this question, we must understand
that models are composed of billions of parameters that are
trained with countless data. The model likely learned lots
of texts that includes ``I love dogs'', ``I love coffee'',
and more. However, for a model dedicated for English,
it is highly unlikely that it learned another phrase that
starts with ``Veni, vidi,'' other than ``Veni, vidi, vici'',
the famous line from Julius Caesar. The model probably
assigned lots of logits and converted them to probability
distribution with softmax. While the first case have somewhat
flat distribution, the second case will likely assign ``vici''
with high probability like $0.98$ or $0.99$. Therefore, the
second model has the higher probability of returning the
expected output. 

Now that we have some intuitions on model outputs, let's
discuss about post-training.

\end{document}


